\section{Conclusion}
\label{sec:conclusion}

The question this paper set out to answer was how \odag{} relates to Formal
Concept Analysis. The answer turned out to be more specific than expected, and
more useful.

\paragraph{They compute the same function.} \fca{}'s derivation operator on
attribute sets \emph{is} \odag{}'s query primitive:
$B' = \cmd{get}(B)$, exactly (Proposition~\ref{prop:dictionary}). An \fca{}
intent \emph{is} an ancestor-set code; an extent \emph{is} a descendant cone;
and the concept lattice of an \odag{} is its Dedekind--MacNeille completion
(Theorem~\ref{thm:dm}). The reflexive ancestor sets are a lossless re-encoding
of the graph, verified on 200 random cases. So the intuition that started the
project --- that a sparse binary code carries a category structure in its
overlaps --- is exactly right, and \fca{} is its mathematics.

\paragraph{They differ in when the function is computed.} \fca{} evaluates it
for every closed set, in advance, and stores the results. \odag{} evaluates it
per query and stores nothing. Everything else follows: the size difference
(linear against up to exponential, measured in Section~\ref{sec:explosion}),
the canonical-form difference (a unique transitive reduction you can hash
against a complete lattice you cannot), the identity difference (stable names
against extent--intent pairs that do not survive an edit), and the
mergeability difference (a CRDT against a structure with no merge at all).

\paragraph{Almost all of \fca{}'s theory is adoptable; almost none of its
representation is.} Section~\ref{sec:learn} found seven imports worth making
--- pattern structures as the correct account of typed values, the scaling
vocabulary and the ordinal kind it reveals as missing, stability and support
as selection criteria, implication bases as a local lint, reduced labelling
for diagrams, the complexity results as guardrails, and attribute exploration
as a protocol --- and every one of them lands in a derived, local, regenerable
layer, or in the documentation, or in the interface. None of them lands in the
stored graph. That is not a coincidence; it is the discipline that made the
list safe to write.

\paragraph{The refusals are the product.} \odag{}'s value to an AI agent is
that its knowledge is canonical, addressable, verifiable and attributable ---
four properties a language model's knowledge lacks, and four properties that
exist only because the system declines to store derived content, declines to
put numbers in identity, declines to let merge fail, and declines to grow past
the largest fragment where ``equal knowledge, equal root'' stays cheap and
semantic. Its value to a decentralised system is the same list read
differently: same meaning gives same hash, so agreement is a comparison and
not a negotiation; merge converges, so consensus is needed only for money;
and there are no floats anywhere, so the structure is unusually ready for the
proof systems nobody has asked for yet.

\medskip

If there is one sentence to carry away, it is the trade in
Section~\ref{sec:identity}, because it explains why two formalisms this
closely related ended up serving different masters:

\begin{keyidea}
\fca{} derives concepts and therefore cannot name them stably.
\odag{} names concepts and therefore cannot derive them.

\fca{} is a discovery instrument. \odag{} is an agreement instrument.
\end{keyidea}

The interesting engineering is in connecting them --- letting \fca{} discover,
in a derived layer over a pinned root, and letting an accountable author
promote what it discovers into a graph that anyone can agree with, cite,
verify and merge. Neither half does that alone. Section~\ref{sec:directions}
is a list of ways to build the joint, and the one worth starting is the oldest
idea in either field: let the machine ask the questions, in minimal
non-redundant order, and let a named author answer them.
